% Generated draft: requires mathematical and source review.
\documentclass[11pt]{article}
\usepackage[margin=25mm]{geometry}
\usepackage{fontspec}
\setmainfont{Latin Modern Roman}
\usepackage{amsmath,amssymb,mathtools,unicode-math}
\setmathfont{Latin Modern Math}
\usepackage{enumitem,needspace,xurl,hyperref}
\setlength{\parindent}{0pt}
\setlength{\parskip}{5pt}
\setlength{\emergencystretch}{3em}
\newcommand{\N}{\mathbb{N}}
\newcommand{\Z}{\mathbb{Z}}
\newcommand{\Q}{\mathbb{Q}}
\newcommand{\R}{\mathbb{R}}
\newcommand{\C}{\mathbb{C}}
\newcommand{\F}{\mathbb{F}}
\newcommand{\A}{\mathbb{A}}
\newcommand{\PP}{\mathbb{P}}
\newcommand{\E}{\mathbb{E}}
\newcommand{\eps}{\varepsilon}
\newcommand{\dd}{\,\mathrm{d}}
\begin{document}
\section*{UnsolvedMath ID: 30002278}
\textbf{Source problem code:} catalogue entry
\subsection{Definitions and mathematical statement}
The supplied record concerns a three-dimensional acoustic perfectly matched layer (PML) system with spatially varying damping coefficients. Its equations, spatial domain, boundary conditions, admissible damping profiles, acoustic energy, and parameterized Lyapunov functional are not supplied. Consequently, the following expresses the requested estimates schematically rather than specifying a complete PDE problem.

Write \(U(t)\) for a solution of the intended system, \(E[U](t)\) for its acoustic energy, and \(\eta_\theta[U](t)\) for the intended Lyapunov functional with parameters \(\theta\). These symbols refer to quantities whose definitions must be recovered from the source; no particular equations or formulas are assumed here. For each admissible system and damping profile, the requested target is to choose admissible parameters \(\theta\) and a constant \(0<C<\infty\), fixed independently of the solution and of time, such that every admissible solution satisfies
\[
0\le E[U](t)\le C\,\eta_\theta[U](t)
\]
and
\[
\eta_\theta[U](t)\le\eta_\theta[U](s)
\qquad\text{whenever }0\le s\le t
\text{ belong to its interval of existence}.
\]
In particular, the argument must control the terms generated by spatial variation of the damping coefficients. The admissible parameter choices, the solution class, and the permitted dependence of \(C\) and \(\theta\) on the damping profiles, domain, and physical coefficients remain unspecified. Uniformity over a class of damping profiles cannot be asserted from the supplied record.

The supplied original question is: ``Can the parameters defining the Lyapunov function \(\eta\) be chosen so that the energy is bounded by a fixed multiple of \(\eta\)?'' The supplied clean statement additionally explicitly requests construction of a nonincreasing functional. Without the functional and its surrounding source discussion, it is unclear whether nonincrease is already established for the intended parameterized family or is itself part of the unresolved task. The source background describes the obstruction using spatially varying conductivity, whereas the clean statement uses spatially varying damping; their precise relationship is not supplied.
\subsection{Short English statement}
For the intended three-dimensional acoustic PML system, choose the parameters of a Lyapunov functional so that it never increases and bounds the acoustic energy by a fixed multiple, including when damping varies in space. The supplied record omits the equations and the energy and functional formulas, so the estimates above describe the target schematically and require source clarification before a complete mathematical statement can be given.
\subsection{Sources}
\begin{itemize}
\item Supplied local UnsolvedMath registry record: permanent ID \(30002278\), source problem code \texttt{catalogue entry}, title \emph{Lyapunov Functions for Three-Dimensional Acoustic PML Systems}, in the Oberwolfach Reports Open Problems collection. The record gives proposed year 2013 and names no proposer.
\item Supplied report citation: \emph{Computational Electromagnetism and Acoustics} (2013), Oberwolfach Reports, DOI \texttt{10.4171/owr/2013/03}, pp.\ 129--237. Source URL: \url{https://doi.org/10.4171/owr/2013/03}. The supplied background describes Oberwolfach Reports as documenting mathematical workshops and their open-problem discussions.
\item Supplied literature source, labeled ``original Oberwolfach report'': \url{https://doi.org/10.4171/owr/2013/03}.
\item Literature assessment date: \texttt{2026-08-21}. This is supplied dated curation, not independently verified. The supplied assessment says that the report identifies an uncontrolled term caused by spatially varying conductivity and that no general resolution was located. The supplied status is ``open''; neither that status nor the literature assessment is independently established here.
\item The registry supplies the statement-status label \texttt{corrected\_verified} and a claim that the clean statement was checked against the official report and preserves its hypotheses, notation, and open target. These are supplied verification claims, not independent verification performed for this draft.
\end{itemize}
\end{document}
